\chapter{Worksheets}

Worksheets are intended as student artifacts. They may be printed, copied into a lab notebook, or reproduced as CSV/Markdown files. Every worksheet should include the date, student name, project folder, and current file path.

\section{Project Brief}
\begin{programbox}{Purpose}
Define a feasible water-quality research question and the minimum viable project.
\end{programbox}
\textbf{Use during:} Week 1, then revise at every decision gate.

\textbf{Required fields:} project title; student name; co-advisers; research question; sample unit; target variable; candidate predictors; project data source; teaching-data fallback; optional uafR role; ML target; expected figures; expected final claim; known limitations; minimum viable product; advanced extension.

\textbf{Common mistakes:} asking a causal question without a causal design; choosing a target that is not measured; including too many systems; treating tentative chemical matches as confirmed.

\textbf{Completion standard:} A reviewer can tell what will be analyzed, why it matters to WiCCED, what data are needed, and what claim will remain out of bounds.

\section{Water-Quality Variable Dictionary}
\textbf{Purpose:} Make every variable interpretable before analysis.

\textbf{Required fields:} column name; plain-language name; unit; variable type; source file; method; sample unit; expected range; missing-value code; detection limit if relevant; target/predictor/excluded status; caveat.

\textbf{Completion standard:} Every feature used in a model appears in the dictionary.

\section{Site and Sample Metadata Sheet}
\textbf{Purpose:} Preserve context that affects interpretation.

\textbf{Required fields:} site ID; site name; latitude/longitude or protected location code; date/time; depth; weather/rainfall note; tidal or hydrologic context if available; sampler; instrument; method; project restrictions; data privacy status.

\textbf{Completion standard:} The sample unit can be reconstructed without relying on memory.

\section{QA/QC Flag Log}
\textbf{Purpose:} Record suspicious values without deleting raw observations.

\textbf{Required fields:} file; row ID; variable; raw value; flag type; rule; reason; action taken; replacement value if any; reviewer; date.

\textbf{Completion standard:} Every changed or excluded value can be traced back to the raw value and decision rule.

\section{Database Evidence Map}
\textbf{Purpose:} Clarify what each source can support.

\textbf{Required fields:} source name; query term; source type; returned field; identifier; what it supports; what it cannot support; citation or URL; planned use in analysis.

\textbf{Completion standard:} No source is used in a claim without a recorded limitation.

\section{GC-MS Input Column Audit}
\textbf{Purpose:} Check whether a GC-MS table can enter uafR.

\textbf{Required fields:} file name; observed columns; required columns; missing columns; renamed columns; duplicate peak keys; suspicious match factors; suspicious areas; decision.

\textbf{Completion standard:} \code{spreadOut()} is run only after required columns and duplicate-key risks are inspected.

\section{Query Settings Log}
\textbf{Purpose:} Make database enrichment reproducible.

\textbf{Required fields:} date; function; compounds; spelling; cache setting; cache directory; throttle; profile/detail; query batch; network status; saved output path; validation output path; source coverage note.

\textbf{Completion standard:} A future student can rerun or inspect the same query without guessing settings.

\section{Trait Evidence Audit}
\textbf{Purpose:} Decide which chemical traits are safe to analyze.

\textbf{Required fields:} compound; trait key; source table; source field; evidence text or ID; extraction rule; confidence; allowed values; analysis use; caveat; include/exclude decision.

\textbf{Completion standard:} Every trait used in a matrix, figure, or model has an evidence row.

\section{ML Feature Inventory}
\textbf{Purpose:} Define target and predictors before modeling.

\textbf{Required fields:} target variable; prediction task; feature name; source; unit/encoding; transformation; available at prediction time; leakage risk; missingness; include/exclude decision.

\textbf{Completion standard:} The feature table is built only from approved inventory rows.

\section{Train/Test Split Plan}
\textbf{Purpose:} Decide what kind of generalization the model is testing.

\textbf{Required fields:} split type; grouping variable; date cutoff if used; training rows; testing rows; reason; leakage risks; baseline model; metrics.

\textbf{Completion standard:} The split design is chosen before model comparison.

\section{Model Card}
\textbf{Purpose:} Summarize model purpose, data, performance, and limitations.

\textbf{Required fields:} model name; target; task type; data version; rows included/excluded; features; split design; algorithm; baseline; metrics; diagnostic plots; limitations; acceptable use; prohibited use; next validation step.

\textbf{Completion standard:} The model can be interpreted without opening the R script first.

\section{Advanced ML Review Gate}
\textbf{Purpose:} Decide whether advanced ML outputs improve the analysis or only add complexity.

\textbf{Use during:} Week 6, after the beginner model card is complete.

\textbf{Required fields:} beginner prerequisite check; advanced script outputs reviewed; time-holdout result; leave-one-site-out result; feature-set sensitivity result; conductivity decision; context-feature decision; top permutation-importance signals; skipped model warnings; final model decision; prohibited claims; next validation need.

\textbf{Common mistakes:} choosing the highest-performing sensitivity model without leakage review; treating permutation importance as causation; ignoring site-holdout weakness; hiding skipped overfit models; using synthetic teaching results as environmental evidence.

\textbf{Completion standard:} A reviewer can see why the preferred model was chosen, why rejected feature sets were rejected, and what claims remain out of bounds.

\section{Figure Planning Worksheet}
\textbf{Purpose:} Plan figures around claims and evidence.

\textbf{Required fields:} figure number; purpose; variables; data source; groups; plot type; units; sample counts; missingness; caption claim; caveat; script path; output path.

\textbf{Completion standard:} Every final figure has a planned claim and caveat.

\section{Figure Caption Checklist}
\textbf{Purpose:} Prevent vague or overstated captions.

\textbf{Required fields:} figure number; variables named; units shown; sample count; split design if model output; missingness noted; color/group definitions; claim strength; limitation sentence.

\textbf{Completion standard:} Captions are understandable without the main text.

\section{Literature and Evidence Synthesis}
\textbf{Purpose:} Connect results to published or official context.

\textbf{Required fields:} citation; system; method; variables; main result; limitation; relevance to project; difference from project; usable sentence.

\textbf{Completion standard:} Literature is used to contextualize, not to claim proof beyond the student's data.

\section{Claim-Evidence-Limitation Memo}
\textbf{Purpose:} Convert results into defensible writing.

\textbf{Required fields:} proposed claim; evidence source; supporting figure/table; limitation; safer wording; follow-up needed.

\textbf{Completion standard:} Every major report claim has a limitation and source.

\section{Peer Review Form}
\textbf{Purpose:} Give useful feedback on project clarity and rigor.

\textbf{Required fields:} reviewer; artifact reviewed; strongest part; unclear part; data concern; model concern; claim concern; reproducibility concern; required revision; optional improvement.

\textbf{Completion standard:} Feedback points to specific files, figures, or statements.

\section{Daily Research Log}
\textbf{Purpose:} Preserve the record of work.

\textbf{Required fields:} date; goal; files opened; commands/scripts run; outputs created; errors; decisions; evidence consulted; interpretation limit; next action.

\textbf{Completion standard:} The log makes the day's work reconstructable.

\section{Weekly Progress Review}
\textbf{Purpose:} Convert daily activity into milestone progress.

\textbf{Required fields:} week; milestone; deliverables completed; files produced; review feedback; changes made; blockers; next week priorities.

\textbf{Completion standard:} The next week starts with a clear action list.

\section{Final Package Audit}
\textbf{Purpose:} Confirm the project can be archived and continued.

\textbf{Required fields:} README present; raw data inventory; processed data; scripts; results; figures; report; presentation; session info; model card; evidence map; continuation note; rerun status; unresolved issues.

\textbf{Completion standard:} A future student can rerun the project or understand why specific steps require restricted data.
